\documentclass[11pt,a4paper]{article}

\usepackage{fontspec}
\setmainfont{DejaVu Serif}
\setsansfont{DejaVu Sans}
\setmonofont{DejaVu Sans Mono}
\usepackage[russian]{babel}
\usepackage{geometry}
\geometry{left=20mm,right=20mm,top=18mm,bottom=20mm}
\usepackage{microtype}
\usepackage{array,tabularx,longtable,booktabs,multirow}
\usepackage{enumitem}
\usepackage{graphicx}
% Renders exported from the Onshape model live next to the STEP files.
\graphicspath{{../hardware/enclosure/renders/}}
\usepackage{xcolor}
\usepackage{hyperref}
\usepackage{fancyhdr}
\usepackage{titlesec}
\usepackage{lastpage}
\usepackage{tcolorbox}
\usepackage{caption}
\usepackage{float}
\usepackage{qrcode}
\usepackage{tikz}
\usetikzlibrary{positioning}

% Build identity is normally injected by `make -C docs build_info`.
% Fallbacks below let the .tex compile standalone in a TeX editor without make.
\IfFileExists{build_info.tex}{\input{build_info}}{%
  \newcommand{\hwversion}{v2.1}%
  \newcommand{\fwversion}{dev}%
  \newcommand{\builddate}{unknown}%
  \newcommand{\buildcommit}{unknown}%
  \newcommand{\docrev}{unknown}%
}

\hypersetup{colorlinks=true,linkcolor=black,urlcolor=blue,pdfauthor={Изготовитель DC-UPS-2CH},pdftitle={Паспорт и руководство по эксплуатации DC-UPS-2CH \fwversion}}
\setlength{\parindent}{0pt}
\setlength{\parskip}{0.45em}
\renewcommand{\arraystretch}{1.2}

\definecolor{accent}{HTML}{1E5B4F}
\definecolor{lightaccent}{HTML}{EAF3F1}
\definecolor{warn}{HTML}{8B1E1E}
\definecolor{lightwarn}{HTML}{F9ECEC}

\titleformat{\section}{\Large\bfseries\color{accent}}{\thesection.}{0.5em}{}
\titleformat{\subsection}{\large\bfseries}{\thesubsection.}{0.5em}{}
\titleformat{\subsubsection}{\normalsize\bfseries}{\thesubsubsection.}{0.5em}{}

\pagestyle{fancy}
\fancyhf{}
\lhead{DC-UPS-2CH}
\rhead{Паспорт и руководство}
\cfoot{Стр. \thepage}
\renewcommand{\headrulewidth}{0.4pt}
\setlength{\headheight}{14pt}

\newcommand{\fieldline}[1]{\textbf{#1}: \rule{0.55\linewidth}{0.4pt}}
\newcommand{\blankfig}[2]{%
  \begin{figure}[H]\centering
  \fbox{\parbox[c][#1][c]{0.90\textwidth}{\centering\sffamily\color{gray}\Large #2\\[0.5em]\normalsize Вставить скриншот / изображение}}
  \end{figure}
}
% \realfig{filename}{caption} — inserts an image full text-width with a caption line.
\newcommand{\realfig}[2]{%
  \begin{figure}[H]\centering
  \includegraphics[width=0.95\textwidth,keepaspectratio]{#1}\\[0.4em]
  {\sffamily\small #2}
  \end{figure}
}
\newcommand{\yesno}{□ да\quad □ нет}

\begin{document}

% -------------------- TITLE --------------------
\begin{titlepage}
\centering
\vspace*{18mm}
{\sffamily\bfseries\fontsize{28}{34}\selectfont DC-UPS-2CH\\[3mm]}
{\sffamily\Large Источник резервного питания сетевого оборудования\\[12mm]}

\begin{tcolorbox}[width=0.85\textwidth,colback=lightaccent,colframe=accent,boxrule=0.8pt,arc=2mm]
\centering
{\Large\bfseries ПАСПОРТ ИЗДЕЛИЯ\\
РУКОВОДСТВО ПО ЭКСПЛУАТАЦИИ\\
СЕРВИСНАЯ ИНСТРУКЦИЯ\\
ГАРАНТИЙНЫЙ ТАЛОН}
\end{tcolorbox}

\vfill
\begin{tabular}{@{}p{55mm}p{90mm}@{}}
Модель изделия: & DC-UPS-2CH \hwversion \\
Версия прошивки: & \fwversion \\
Редакция документа: & \builddate\ / \docrev \\
Серийный номер: & \rule{70mm}{0.4pt} \\
Дата изготовления: & \rule{70mm}{0.4pt} \\
Изготовитель / сборщик: & \rule{70mm}{0.4pt} \\
Контакт: & \rule{70mm}{0.4pt} \\
\end{tabular}

\vfill
{\small Документ предназначен для владельца изделия и для выполнения регламентного обслуживания.\\
Хранить вместе с устройством. Актуализировано для прошивки \fwversion\ (сборка \buildcommit, \builddate).}
\end{titlepage}

\tableofcontents
\newpage

% -------------------- 1 --------------------
\section{Общие сведения}

\subsection{Назначение}
DC-UPS-2CH представляет собой двухканальный источник резервного питания постоянного тока для сетевого и слаботочного оборудования. Устройство предназначено для поддержания работы Wi-Fi роутера, ONT/оптического терминала, модема или другой совместимой нагрузки при пропадании сетевого питания 230~В.

Устройство сочетает сетевой источник 24~В, зарядный преобразователь, герметичный аккумулятор 12~В, два независимых выходных DC/DC-канала и управляющий контроллер ESP32. Контроллер отслеживает состояние сети и аккумулятора, управляет двумя выходами, ведёт журнал событий, выполняет автоматическое восстановление связи и предоставляет веб-интерфейс.

\subsection{Основные функции}
\begin{itemize}[leftmargin=7mm]
  \item автоматическая работа от сети и переход на аккумулятор без ручного переключения;
  \item поддерживающий заряд аккумулятора с ограничением тока;
  \item два независимых управляемых выходных канала ROUTER и ONT;
  \item защита аккумулятора от глубокого разряда (LVD);
  \item аварийный deep sleep контроллера после отключения нагрузки;
  \item периодическая проверка возврата сети во время аварийного deep sleep;
  \item отдельный режим хранения (Shelf Sleep): оба выхода и Wi-Fi выключены, пробуждение только физической кнопкой;
  \item контроль наличия Wi-Fi и доступности интернета;
  \item автоматический power-cycle выходных каналов при неисправности связи;
  \item локальный веб-интерфейс;
  \item mDNS-адрес \texttt{http://dc-ups.local/};
  \item сервисная точка доступа \texttt{DC-UPS-Setup};
  \item уведомления через сервис ntfy;
  \item read-only команды через тот же ntfy topic: \texttt{!ups}, \texttt{!ups status}, \texttt{!ups ping};
  \item журнал событий в оперативной памяти;
  \item локальная тактовая кнопка: статус, setup AP и режим хранения;
  \item световая индикация состояния сети и аккумулятора.
\end{itemize}

\begin{tcolorbox}[colback=lightwarn,colframe=warn,title=Важно]
Внутри устройства присутствует опасное сетевое напряжение 230~В. Разборка, регулировка и замена внутренних компонентов выполняются только после отключения сетевого питания и аккумулятора.
\end{tcolorbox}

% -------------------- 2 --------------------
\section{Технические характеристики}

\subsection{Основные параметры}
\begin{longtable}{@{}p{70mm}p{80mm}@{}}
\toprule
\textbf{Параметр} & \textbf{Значение} \\
\midrule
\endhead
Тип изделия & Двухканальный DC-UPS для слаботочного оборудования \\
Сетевое питание & 230~В AC, 50~Гц \\
Промежуточная шина & 24~В DC от внутреннего AC/DC-модуля \\
Номинал аккумулятора & 12~В, 7~А\,ч, герметичный AGM/SLA \\
Установленная модель АКБ & \rule{60mm}{0.4pt} \\
Напряжение поддерживающего заряда & 13,8~В на клеммах АКБ (настроечное значение) \\
Ограничение тока зарядного канала & 1,5~А \\
Зарядный DC/DC-модуль & XL4016 CC/CV или установленный эквивалент \\
Количество выходов & 2 независимых канала \\
Выходные DC/DC-модули & 2\,$\times$\,XL6009 buck-boost или установленный эквивалент \\
Управляющий контроллер & ESP32 \\
Версия прошивки на момент выпуска документа & \fwversion \\
Контроль АКБ & ADC, GPIO35, делитель 100~кОм / 10~кОм \\
Контроль 24~В & ADC, GPIO34, делитель 100~кОм / 10~кОм \\
Канал ROUTER & GPIO32 $\rightarrow$ NPN 2N2222 $\rightarrow$ P-MOSFET \\
Канал ONT & GPIO25 $\rightarrow$ NPN 2N2222 $\rightarrow$ P-MOSFET \\
Сервисная кнопка & GPIO14, замыкание на GND, INPUT\_PULLUP; RTC GPIO / EXT0 wake в режиме хранения \\
LED сети & GPIO13 \\
LED АКБ/аварии & GPIO12 \\
Подтяжка затвора P-MOSFET & 10~кОм Gate--Source (по текущей схеме) \\
Модель P-MOSFET & IRF9Zxx / фактически установленный: \rule{35mm}{0.4pt} \\
Предохранитель АКБ & рекомендуемый 5~А, установить непосредственно у плюсовой клеммы \\
\bottomrule
\end{longtable}

\subsection{Параметры выходов конкретного экземпляра}
\begin{tabularx}{\textwidth}{@{}lXccc@{}}
\toprule
\textbf{Канал} & \textbf{Назначение / подключённое устройство} & \textbf{U, В} & \textbf{I nom, А} & \textbf{I max, А} \\
\midrule
ROUTER & \rule{45mm}{0.4pt} & \rule{12mm}{0.4pt} & \rule{12mm}{0.4pt} & \rule{12mm}{0.4pt} \\
ONT & \rule{45mm}{0.4pt} & \rule{12mm}{0.4pt} & \rule{12mm}{0.4pt} & \rule{12mm}{0.4pt} \\
\bottomrule
\end{tabularx}

\begin{tcolorbox}[colback=lightaccent,colframe=accent,title=Примечание по выходным токам]
Максимальный допустимый ток каждого выхода определяется фактически установленным DC/DC-модулем, охлаждением, проводкой, разъёмом и P-MOSFET. Перед передачей изделия владельцу значения должны быть установлены по результатам нагрузочного испытания и вписаны в таблицу и на шильдик. Не следует указывать номинал только по рекламному ``максимальному току'' модуля XL6009.
\end{tcolorbox}

\subsection{Пороговые значения защиты}
\begin{tabularx}{\textwidth}{@{}XlX@{}}
\toprule
\textbf{Функция} & \textbf{Значение по умолчанию} & \textbf{Назначение} \\
\midrule
Предупреждение АКБ & 11,5~В & Уведомление о приближении разряда \\
LVD / отключение нагрузки & 11,0~В & Отключение ROUTER и ONT при отсутствии сети \\
Восстановление после LVD & 12,5~В & Разрешение повторного включения после восстановления АКБ \\
Дополнительный аварийный порог & 10,8~В & Второй уровень контроля низкого напряжения \\
Задержка перед deep sleep & 30~с & Даёт время завершить уведомление и стабилизировать состояние \\
Период пробуждения в аварийном deep sleep & 30~с & Проверка возврата 24~В и состояния АКБ \\
Shelf Sleep / режим хранения & без периодического таймера & Пробуждение только кнопкой GPIO14 (EXT0) \\
``Сеть есть'' & выше 15,0~В на контролируемой 24-В линии & Верхний порог гистерезиса \\
``Сети нет'' & ниже 10,0~В на контролируемой 24-В линии & Нижний порог гистерезиса \\
Антидребезг сети & 3~с & Исключение ложных переключений \\
\bottomrule
\end{tabularx}

Все перечисленные пороги, кроме аппаратных номиналов, могут быть изменены в веб-интерфейсе. Значения в данном документе соответствуют настройкам прошивки по умолчанию.

% -------------------- 3 --------------------
\section{Структура и принцип действия}

\subsection{Упрощённая структурная схема}
\begin{center}
\begin{tikzpicture}[scale=0.78,transform shape,>=latex, node distance=7mm and 7mm, every node/.style={font=\sffamily\small}]
\node[draw,rounded corners,minimum width=28mm,minimum height=10mm] (ac) {230 В AC};
\node[draw,rounded corners,right=of ac,minimum width=32mm,minimum height=10mm] (psu) {AC/DC 24 В};
\node[draw,rounded corners,right=of psu,minimum width=34mm,minimum height=10mm] (chg) {XL4016 CC/CV};
\node[draw,rounded corners,right=of chg,minimum width=30mm,minimum height=10mm] (bus) {13,8 В / АКБ};
\node[draw,rounded corners,below=12mm of bus,minimum width=30mm,minimum height=10mm] (bat) {AGM 12 В 7 А·ч};
\node[draw,rounded corners,below left=18mm and 15mm of bus,minimum width=36mm,minimum height=10mm] (r) {P-MOSF 1 + XL6009};
\node[draw,rounded corners,below right=18mm and 15mm of bus,minimum width=36mm,minimum height=10mm] (o) {P-MOSF 2 + XL6009};
\node[draw,rounded corners,below=13mm of r,minimum width=28mm,minimum height=9mm] (rr) {ROUTER};
\node[draw,rounded corners,below=13mm of o,minimum width=28mm,minimum height=9mm] (oo) {ONT / AUX};
\node[draw,rounded corners,below=35mm of bus,minimum width=34mm,minimum height=10mm,fill=lightaccent] (esp) {ESP32};
\draw[->] (ac)--(psu); \draw[->] (psu)--(chg); \draw[->] (chg)--(bus); \draw[<->] (bus)--(bat);
\draw[->] (bus)--(r); \draw[->] (bus)--(o); \draw[->] (r)--(rr); \draw[->] (o)--(oo);
\draw[->,dashed] (esp)--node[left]{GPIO32}(r); \draw[->,dashed] (esp)--node[right]{GPIO25}(o);
\draw[->,dashed] (bus)--node[right]{ADC}(esp);
\end{tikzpicture}
\end{center}

\subsection{Штатный режим}
При наличии сетевого питания внутренний AC/DC-модуль формирует 24~В. Зарядный модуль XL4016 понижает напряжение и поддерживает на аккумуляторной шине установленное напряжение 13,8~В с ограничением суммарного тока до 1,5~А. Нагрузка и заряд аккумулятора питаются от общей шины после развязывающего диода согласно фактической сборке.

\subsection{Работа от аккумулятора}
После пропадания 24-В линии аккумулятор автоматически становится источником питания общей шины. Отдельного реле переключения не требуется. Контроллер фиксирует пропадание сети, оставляет разрешённые каналы включёнными и отправляет уведомление при наличии связи.

\subsection{Защита от глубокого разряда}
При работе без сети контроллер сравнивает напряжение АКБ с установленными порогами. При достижении LVD оба выхода отключаются. Повторное включение после срабатывания LVD не разрешается до достижения порога восстановления либо возврата сети.

После LVD и снятия нагрузки контроллер через установленную задержку переходит в \textbf{аварийный deep sleep}, не ожидая повторного падения напряжения после естественного ``отскока'' разгруженного аккумулятора. В этом режиме ESP32 периодически просыпается по таймеру только для измерения АКБ и 24-В линии. Полный запуск выполняется при возврате сети либо при восстановлении АКБ до порога разрешённого запуска.

\subsection{Два режима сна}
\textbf{Аварийный deep sleep} запускается автоматически при низком АКБ и отсутствии сети. Оба выхода выключены, Wi-Fi выключен, контроллер просыпается через заданный интервал (по умолчанию 30~с), быстро измеряет АКБ и линию 24~В и снова засыпает, если условия восстановления не выполнены.

\textbf{Shelf Sleep / режим хранения} запускается только длительным удержанием сервисной кнопки. В этом режиме оба выхода, Wi-Fi, mDNS и setup AP выключаются, таймер пробуждения не используется. Единственный источник пробуждения --- физическая кнопка GPIO14 через EXT0. Возврат сетевого питания 230~В сам по себе \textbf{не выводит устройство из режима хранения}. Для включения необходимо нажать кнопку.

% -------------------- 4 --------------------
\section{Органы управления и индикация}

\subsection{Светодиоды}
\begin{tabularx}{\textwidth}{@{}lX@{}}
\toprule
\textbf{Индикатор} & \textbf{Назначение} \\
\midrule
Зелёный, GPIO13 & Индикация наличия внешнего питания / состояния сети \\
Красный, GPIO12 & Индикация состояния АКБ, предупреждения и аварийного режима \\
\bottomrule
\end{tabularx}

\subsection{Сервисная кнопка GPIO14}
Кнопка подключается между GPIO14 и GND. В обычной работе используется внутренняя подтяжка \texttt{INPUT\_PULLUP}. GPIO14 выбран также как RTC GPIO для пробуждения из режима хранения.

\begin{tabularx}{\textwidth}{@{}p{38mm}X@{}}
\toprule
\textbf{Действие} & \textbf{Функция} \\
\midrule
Одно короткое нажатие & Отправить полный текущий статус устройства в настроенный ntfy topic. После успешной отправки зелёный LED мигает 2 раза; если Wi-Fi/ntfy недоступны, красный LED мигает 3 раза. \\
Двойное нажатие & Запустить \texttt{DC-UPS-Setup} примерно на 15 минут. Подтверждение --- краткое совместное мигание обоих LED. \\
Удержание 10~с & Перевести устройство в Shelf Sleep / режим хранения. Перед выключением красный LED мигает 3 раза; затем ROUTER и ONT выключаются, Wi-Fi отключается и ESP32 засыпает без таймера. \\
Нажатие в режиме хранения & Пробудить устройство. После пробуждения удерживаемое нажатие игнорируется до отпускания кнопки, затем выполняется штатный запуск. \\
\bottomrule
\end{tabularx}

\begin{tcolorbox}[colback=lightwarn,colframe=warn,title=Важно: режим хранения]
В Shelf Sleep устройство \textbf{не просыпается от появления 230~В}. Оно намеренно остаётся выключенным до физического нажатия кнопки. Это режим длительного хранения / перевозки, а не аварийный режим разряда АКБ.
\end{tcolorbox}

\subsection{Ручные выключатели выходов}
Если на корпусе установлены отдельные механические выключатели ROUTER и ONT, они считаются локальными аппаратными выключателями соответствующего выхода. Их положение имеет приоритет на уровне физического разрыва цепи: программное включение не восстановит питание при разомкнутом механическом выключателе.

% -------------------- 5 --------------------
\section{Веб-интерфейс и сетевые функции}

\subsection{Доступ}
При штатном подключении к локальной сети панель доступна по адресу:
\begin{center}
\texttt{http://dc-ups.local/}
\end{center}
Также панель доступна по IP-адресу, назначенному роутером. При каждом успешном подключении ESP32 к Wi-Fi прошивка может отправлять актуальный IP через ntfy.

\subsection{Сервисная точка доступа}
SSID по умолчанию: \texttt{DC-UPS-Setup}.\\
Пароль по умолчанию: \texttt{dc-ups-setup}.\\
Параметры могут быть изменены владельцем через веб-интерфейс.

\begin{tcolorbox}[colback=lightaccent,colframe=accent,title=Рекомендация]
Пароль setup AP и пароль администратора веб-панели не следует печатать на наружной наклейке устройства. Их лучше передать владельцу отдельно и изменить после ввода в эксплуатацию.
\end{tcolorbox}

\subsection{Двусторонний обмен через ntfy}
Помимо исходящих уведомлений прошивка периодически опрашивает тот же ntfy topic (примерно один раз в 15~с) и принимает только безопасные read-only команды состояния:
\begin{center}
\texttt{!ups} \qquad \texttt{!ups status} \qquad \texttt{!ups ping}
\end{center}
Команды \texttt{!ups} и \texttt{!ups status} вызывают отправку полного статуса. Команда \texttt{!ups ping} возвращает сообщение ``PONG'' и тот же набор параметров. В ответ включаются напряжение АКБ, напряжение 24~В, состояние сети, ROUTER/ONT, интернет, recovery, uptime, число отключений, RSSI, IP и mDNS-адрес.

Удалённые команды включения/выключения выходов через ntfy \textbf{не реализованы}. Это сделано намеренно: topic используется только как канал чтения статуса. Команды работают только когда ESP32 бодрствует, подключена к Wi-Fi и имеет доступ к ntfy. В Shelf Sleep и во время аварийного сна Wi-Fi выключен, поэтому команды недоступны.

\begin{tcolorbox}[colback=lightaccent,colframe=accent,title=Безопасность ntfy]
Имя topic фактически является секретным адресом. Использовать длинное, трудноугадываемое имя и не печатать его на наружной наклейке устройства.
\end{tcolorbox}

\subsection{Доступные функции веб-панели}
\begin{itemize}[leftmargin=7mm]
  \item напряжение АКБ и 24-В линии;
  \item наличие сети;
  \item отдельное состояние ROUTER и ONT;
  \item режимы ON / OFF / AUTO каждого канала;
  \item отдельный restart ROUTER и restart ONT;
  \item совместный restart обоих выходов;
  \item состояние Wi-Fi и RSSI;
  \item состояние проверки интернета и autorecovery;
  \item состояние аварийного deep sleep и его параметры;
  \item справка по Shelf Sleep и сервисной кнопке;
  \item калибровка ADC по мультиметру;
  \item настройка Wi-Fi, ntfy и защитных порогов;
  \item справка по входящим ntfy-командам;
  \item журнал событий;
  \item ручная перезагрузка ESP32.
\end{itemize}

% -------------------- 6 --------------------
\section{Алгоритм автоматического восстановления связи}

Прошивка различает потерю Wi-Fi и потерю доступа в интернет. Типовая логика:
\begin{enumerate}[leftmargin=8mm]
  \item Если ESP32 не подключена к Wi-Fi после допустимого времени ожидания, выполняется попытка восстановить соединение.
  \item При устойчивой неисправности Wi-Fi первым перезапускается канал ROUTER.
  \item Если Wi-Fi присутствует, но интернет недоступен, первым перезапускается канал ONT.
  \item При продолжении неисправности допускается следующий цикл восстановления, включая совместный power-cycle обоих каналов.
  \item Количество автоматических power-cycle ограничено, чтобы исключить бесконечное переключение оборудования.
  \item После длительного периода стабильной связи счётчики autorecovery сбрасываются.
  \item При низком АКБ автоматические перезапуски блокируются, чтобы не расходовать остаточный заряд.
\end{enumerate}

% -------------------- 7 --------------------
\section{Установка и ввод в эксплуатацию}

\subsection{Предварительная проверка}
Перед первым включением или после обслуживания:
\begin{enumerate}[leftmargin=8mm]
  \item Проверить полярность аккумулятора.
  \item Проверить наличие и номинал предохранителя АКБ.
  \item Проверить отсутствие незакреплённых проводов и токоведущих частей.
  \item Мультиметром выставить и проверить выходное напряжение каждого XL6009 \textbf{до подключения нагрузки}.
  \item Убедиться, что полярность DC-разъёмов соответствует подключаемому оборудованию.
  \item Проверить, что напряжение ROUTER и ONT записано в паспорт и на шильдик.
  \item Подать 230~В и проконтролировать напряжение на клеммах АКБ.
  \item Проверить переход на аккумулятор кратковременным отключением 230~В.
\end{enumerate}

\subsection{Контрольное испытание перед передачей}
\begin{tabularx}{\textwidth}{@{}XcX@{}}
\toprule
\textbf{Проверка} & \textbf{Результат} & \textbf{Примечание} \\
\midrule
Напряжение заряда АКБ 13,8~В & \yesno & \rule{40mm}{0.4pt} \\
ROUTER ON/OFF/AUTO & \yesno & \rule{40mm}{0.4pt} \\
ONT ON/OFF/AUTO & \yesno & \rule{40mm}{0.4pt} \\
Restart ROUTER из веб-панели & \yesno & \rule{40mm}{0.4pt} \\
Restart ONT из веб-панели & \yesno & \rule{40mm}{0.4pt} \\
Переход сеть $\rightarrow$ АКБ & \yesno & \rule{40mm}{0.4pt} \\
LVD & \yesno & \rule{40mm}{0.4pt} \\
Аварийный deep sleep и периодические wake-up & \yesno & \rule{40mm}{0.4pt} \\
Возврат сети после аварийного deep sleep & \yesno & \rule{40mm}{0.4pt} \\
Shelf Sleep: вход удержанием 10~с & \yesno & \rule{40mm}{0.4pt} \\
Shelf Sleep: wake только кнопкой & \yesno & \rule{40mm}{0.4pt} \\
Кнопка: 1x status / 2x setup AP & \yesno & \rule{40mm}{0.4pt} \\
Wi-Fi / mDNS & \yesno & \rule{40mm}{0.4pt} \\
ntfy: исходящие уведомления & \yesno & \rule{40mm}{0.4pt} \\
ntfy: \texttt{!ups status} / \texttt{!ups ping} & \yesno & \rule{40mm}{0.4pt} \\
Нагрузочное испытание обоих каналов & \yesno & \rule{40mm}{0.4pt} \\
\bottomrule
\end{tabularx}

% -------------------- 8 --------------------
\section{Разборка и замена аккумулятора}

\begin{tcolorbox}[colback=lightwarn,colframe=warn,title=Опасное напряжение]
Перед снятием крышки обязательно отключить устройство от 230~В, отсоединить сетевой шнур, отключить нагрузку и разомкнуть цепь аккумулятора. Даже при отключённой сети аккумулятор остаётся источником энергии и способен выдавать высокий ток короткого замыкания.
\end{tcolorbox}

\subsection{Порядок разборки}
\begin{enumerate}[leftmargin=8mm]
  \item Отключить 230~В и убедиться, что сетевой ввод физически отсоединён.
  \item Выключить внешние нагрузки ROUTER и ONT.
  \item Снять крышку согласно рисункам в разделе~\ref{sec:figures}.
  \item Отключить предохранитель / разъём аккумулятора.
  \item Первой отсоединить клемму, указанную в сервисной схеме конкретной сборки; не допускать контакта инструмента между ``+'' и корпусом/``-''.
  \item Отсоединить вторую клемму и извлечь аккумулятор.
  \item Установить новый AGM/SLA аккумулятор 12~В 7~А\,ч совместимого размера и типа клемм.
  \item Подключить полярность строго по маркировке: красный провод к ``+'', чёрный к ``-''.
  \item Вернуть предохранитель, проверить напряжение мультиметром и собрать корпус.
  \item Выполнить контрольный переход сеть $\rightarrow$ АКБ.
\end{enumerate}

\subsection{Требования к аккумулятору-замене}
\begin{itemize}[leftmargin=7mm]
  \item номинальное напряжение 12~В;
  \item технология AGM/SLA (герметичный свинцово-кислотный);
  \item номинальная ёмкость 7~А\,ч либо другая, допускаемая корпусом и зарядным режимом;
  \item совместимые габариты, расположение клемм и тип наконечников;
  \item отсутствие вздутия, трещин и следов утечки.
\end{itemize}

\subsection{Сервисные изображения и 3D-модель}\label{sec:figures}
Сборочные виды экспортированы из общедоступной модели Onshape (см. подраздел о 3D-модели в README проекта). Полные STEP-файлы лежат в каталоге репозитория \texttt{hardware/enclosure/step/}.

\realfig{exploded-front.png}{Рисунок 1. Разнесённый вид со стороны лицевой панели: корпус из алюминиевого профиля, аккумулятор AGM~12\,В\,7\,А·ч установлен внутри, лицевая панель с входным разъёмом IEC~C14, силовым выключателем и двумя гнёздами DC-выходов вынесена слева.}

\realfig{exploded-rear.png}{Рисунок 2. Разнесённый вид со стороны задней панели: корпус с уже установленными органами управления, задняя крышка снята, показана двухъярусная плата с зарядным модулем XL4016, двумя DC/DC XL6009 и материнской платой ESP32.}

\blankfig{55mm}{Рисунок 3. Винты / защёлки крышки и направление разборки}
\blankfig{55mm}{Рисунок 4. Расположение предохранителя и клемм АКБ}
\blankfig{55mm}{Рисунок 5. Извлечение и установка аккумулятора}
\blankfig{55mm}{Рисунок 6. Платы XL4016, ESP32, два XL6009 и силовые каналы крупным планом}

% -------------------- 9 --------------------
\section{Техническое обслуживание}

Рекомендуемая периодичность проверки: один раз в 3--6 месяцев, а также после длительного отключения сети или заметного сокращения автономности.

\begin{itemize}[leftmargin=7mm]
  \item визуальный осмотр корпуса и кабелей;
  \item проверка отсутствия нагрева, запаха, потемнений и деформации компонентов;
  \item проверка затяжки клемм при отключённом питании;
  \item проверка напряжения заряда на клеммах АКБ;
  \item проверка напряжения обоих выходов;
  \item тест работы от АКБ;
  \item проверка LVD и восстановления;
  \item проверка состояния аккумулятора и даты его установки;
  \item проверка журнала событий и исходящих ntfy-уведомлений;
  \item проверка ответа на \texttt{!ups status};
  \item проверка входа в Shelf Sleep и пробуждения кнопкой.
\end{itemize}

\subsection{Ожидаемый срок службы АКБ}
Срок службы герметичного свинцово-кислотного аккумулятора зависит от температуры, глубины и частоты разрядов и качества поддерживающего заряда. Аккумулятор является расходным элементом и подлежит замене при существенном снижении фактической автономности, деформации корпуса или нестабильном напряжении под нагрузкой.

% -------------------- 10 --------------------
\section{Возможные неисправности}
\begin{longtable}{@{}p{38mm}p{48mm}p{74mm}@{}}
\toprule
\textbf{Признак} & \textbf{Возможная причина} & \textbf{Действия} \\
\midrule
\endhead
Нет питания на обоих выходах & Режим OFF, LVD, отсутствует АКБ, предохранитель, неисправность общей шины & Проверить веб-панель, напряжение АКБ, предохранитель и клеммы \\
Нет ROUTER & Выключен канал 1, механический выключатель, XL6009 канала 1, силовой ключ & Проверить режим ROUTER и напряжение до/после соответствующего DC/DC \\
Нет ONT & Выключен канал 2, механический выключатель, XL6009 канала 2, силовой ключ & Проверить режим ONT и напряжение до/после соответствующего DC/DC \\
Нет веб-панели & ESP не подключилась к Wi-Fi либо роутер отключён & Использовать сервисную кнопку, setup AP и Serial при обслуживании \\
Wi-Fi есть, интернета нет & Сбой ONT/провайдера/WAN & Проверить журнал; устройство может выполнить автоматический restart ONT \\
Быстро разряжается АКБ & Старение АКБ или увеличившаяся нагрузка & Измерить ток/мощность нагрузки и выполнить тест ёмкости \\
Система ушла в deep sleep & Низкий АКБ после LVD при отсутствии сети & Восстановить 230~В и дождаться ближайшей периодической проверки \\
Устройство не включилось после подачи 230~В & Оно находится в Shelf Sleep / режиме хранения & Нажать сервисную кнопку; сеть 230~В намеренно не является источником wake-up в этом режиме \\
Кнопка дала 3 красных мигания & Не удалось отправить статус в ntfy & Проверить Wi-Fi, интернет и настройку topic; само устройство продолжает работать \\
Команда \texttt{!ups status} не отвечает & Устройство спит, нет Wi-Fi/интернета или неверный topic & Разбудить устройство при необходимости, проверить Wi-Fi и topic; команды опрашиваются примерно раз в 15~с \\
Выходное напряжение проседает & Перегрузка XL6009, плохой контакт, силовой ключ или проводка & Снять нагрузку, измерить напряжение по узлам, проверить ток и нагрев \\
\bottomrule
\end{longtable}

% -------------------- 11 --------------------
\section{Комплект поставки}
\begin{tabularx}{\textwidth}{@{}clXc@{}}
\toprule
№ & \textbf{Наименование} & \textbf{Примечание} & \textbf{Кол.} \\
\midrule
1 & DC-UPS-2CH & Основной блок & 1 \\
2 & АКБ 12~В 7~А\,ч AGM/SLA & Установлен в корпус & 1 \\
3 & Паспорт и руководство & Настоящий документ & 1 \\
4 & Гарантийный талон & В составе документа & 1 \\
5 & Кабели питания выходов & \rule{55mm}{0.4pt} & \rule{10mm}{0.4pt} \\
6 & Прочее & \rule{55mm}{0.4pt} & \rule{10mm}{0.4pt} \\
\bottomrule
\end{tabularx}

% -------------------- 12 --------------------
\section{Свидетельство о приёмке}
Изделие DC-UPS-2CH \hwversion, серийный номер \rule{50mm}{0.4pt}, собрано, проверено и признано годным к эксплуатации при соблюдении настоящего руководства.

\vspace{5mm}
\begin{tabular}{@{}p{55mm}p{90mm}@{}}
Дата изготовления: & \rule{65mm}{0.4pt} \\
Дата проверки: & \rule{65mm}{0.4pt} \\
Проверил: & \rule{65mm}{0.4pt} \\
Подпись: & \rule{65mm}{0.4pt} \\
\end{tabular}

% -------------------- 13 --------------------
\newpage
\section{Гарантийный талон}

\begin{tcolorbox}[colback=white,colframe=accent,boxrule=1pt,arc=2mm]
\begin{tabular}{@{}p{55mm}p{90mm}@{}}
Изделие: & DC-UPS-2CH \hwversion \\
Серийный номер: & \rule{65mm}{0.4pt} \\
Дата передачи: & \rule{65mm}{0.4pt} \\
Изготовитель / сборщик: & \rule{65mm}{0.4pt} \\
Владелец: & \rule{65mm}{0.4pt} \\
Гарантийный срок: & \rule{30mm}{0.4pt} месяцев \\
\end{tabular}
\end{tcolorbox}

\subsection*{Гарантия распространяется}
На дефекты сборки и отказ внутренних узлов изделия при соблюдении условий эксплуатации, если иное письменно не оговорено при передаче.

\subsection*{Гарантия не распространяется}
\begin{itemize}[leftmargin=7mm]
  \item на естественное старение и потерю ёмкости аккумулятора;
  \item на повреждения вследствие переполюсовки, короткого замыкания или подключения нагрузки с неверным напряжением;
  \item на перегрузку выходов сверх записанных в паспорте значений;
  \item на попадание воды, конденсат, пожар, грозовые и иные внешние воздействия;
  \item на механические повреждения;
  \item на последствия самостоятельной переделки силовой части или изменения настроек, выходящих за согласованные пределы;
  \item на неисправности внешнего роутера, ONT, сети провайдера и подключённых кабелей.
\end{itemize}

\subsection*{Отметки о ремонте}
\begin{tabularx}{\textwidth}{@{}p{25mm}Xp{35mm}@{}}
\toprule
\textbf{Дата} & \textbf{Выполненные работы} & \textbf{Подпись} \\
\midrule
\rule{0pt}{10mm} & & \\
\midrule
\rule{0pt}{10mm} & & \\
\midrule
\rule{0pt}{10mm} & & \\
\bottomrule
\end{tabularx}

\vfill
\begin{tabular}{@{}p{72mm}p{72mm}@{}}
Подпись изготовителя: \rule{35mm}{0.4pt} & Подпись владельца: \rule{35mm}{0.4pt} \\
\end{tabular}

% -------------------- 14 --------------------
\newpage
\section{Краткая памятка владельцу}
\begin{enumerate}[leftmargin=8mm]
  \item В штатном режиме устройство не требует ручного переключения между сетью и АКБ.
  \item При отключении 230~В работа продолжается от аккумулятора до достижения порога защиты.
  \item После LVD устройство отключает нагрузку и переходит в аварийный deep sleep для сохранения аккумулятора.
  \item После возврата сети устройство автоматически выходит из \textbf{аварийного} deep sleep, включает разрешённые выходы и восстанавливает Wi-Fi.
  \item Панель: \texttt{http://dc-ups.local/}.
  \item Кнопка: 1 короткое -- отправить статус в ntfy; 2 коротких -- setup AP на 15 минут; удержание 10~с -- Shelf Sleep.
  \item В Shelf Sleep устройство просыпается \textbf{только кнопкой}; подача 230~В сама по себе его не включает.
  \item Запрос статуса через ntfy: \texttt{!ups}, \texttt{!ups status} или \texttt{!ups ping}.
  \item Перед подключением нового оборудования проверить требуемое напряжение и полярность.
  \item Не вскрывать корпус при подключённой сети 230~В или подключённом аккумуляторе.
\end{enumerate}

\begin{center}
\vfill
\qrcode[height=28mm]{http://dc-ups.local/}\\[2mm]
{\small Быстрый переход к локальной панели (работает из той же локальной сети при поддержке mDNS).}
\end{center}

\end{document}
